% Adapted from the newer reference talk's aligned-block CNOT picture.
\begin{frame}{Background: Homomorphic CNOT}
  \centering
  \begin{tikzpicture}[
      >={Latex[length=2.6mm]},
      qubit/.style={rounded corners=1.5pt,minimum width=0.56cm,
        minimum height=0.58cm,inner sep=0pt,thick},
      data/.style={qubit,draw=CVQCTeal,fill=CVQCPale},
      ztrap/.style={qubit,draw=CVQCGold,fill=CVQCGoldPale},
      xtrap/.style={qubit,draw=CVQCInk,fill=CVQCFog},
      hidden/.style={qubit,draw=CVQCMuted,fill=CVQCFog}
    ]
    \path[use as bounding box] (-6.6,-3.15) rectangle (6.6,3.05);

    \node[font=\large,text=CVQCInk] at (0.8,2.1) {Same permutation};
    \node[anchor=east,font=\large,text=CVQCInk] at (-4.05,0.75)
      {$\operatorname{Enc}(\ket{\psi})$};
    \node[anchor=east,font=\large,text=CVQCInk] at (-4.05,-0.75)
      {$\operatorname{Enc}(\ket{\phi})$};

    % Four data positions and four of each trap type; physical counts schematic.
    % The revealed colors are for the audience, not information given to the server.
    \foreach \i/\kind in {0/ztrap,1/data,2/xtrap,3/ztrap,4/data,5/xtrap,
        6/xtrap,7/ztrap,8/data,9/ztrap,10/data,11/xtrap} {
      \pgfmathsetmacro{\x}{-3.6+0.8*\i}
      \only<1>{
        \node[hidden] at (\x,0.75) {};
        \node[hidden] at (\x,-0.75) {};
      }
      \only<2>{
        \node[\kind] at (\x,0.75) {};
        \node[\kind] at (\x,-0.75) {};
      }
      \draw[thick,draw=CVQCInk] (\x,0.75) -- (\x,-0.75);
      \fill[CVQCInk] (\x,0.75) circle (2pt);
      \draw[thick,draw=CVQCInk,fill=white] (\x,-0.75) circle (0.14);
      \draw[thick,draw=CVQCInk]
        ({\x-0.14},-0.75) -- ({\x+0.14},-0.75)
        (\x,-0.89) -- (\x,-0.61);
    }

    \draw[decorate,decoration={brace,mirror,amplitude=5pt},
      thick,draw=CVQCMuted] (-3.9,-1.25) -- (5.5,-1.25);
    \node[font=\large,text=CVQCInk] at (0.8,-1.95) {Transversal CNOT};
  \end{tikzpicture}

  \note<1>{
    \footnotesize
    Take two authenticated logical qubits. The first row is the control and
    the second is the target. Both blocks use the same hidden permutation,
    with independent one-time pads. Apply a physical CNOT from each position
    in the first ciphertext to the corresponding position in the second.
    There are 3n physical CNOTs, where n is the underlying code block length;
    the twelve drawn pairs are schematic. The prover can do this without
    learning which positions contain encoded data or either type of trap.
    The verifier updates its pad keys. If the old keys are
    (x1,z1) and (x2,z2), the new keys are
    (x1,z1 XOR z2) and (x1 XOR x2,z2).
  }
  \note<2>{
    \footnotesize
    Reveal the colors only to explain why the operation works. The prover
    never learns them. Matching permutations align data with data, zero traps
    with zero traps, and plus traps with plus traps. The CSS code implements
    the logical CNOT transversally. CNOT preserves each pair of zero traps
    and each pair of plus traps, so the traps remain valid under the updated
    pad keys. The two logical output registers may be entangled; the input
    labels on the left do not assert product-form outputs. This explains
    honest correctness; security against deviations still uses authentication.
    Other gates, especially the non-Clifford resource needed for our
    computation, require additional procedures rather than just transversal
    gates. Figure adapted from the newer reference talk's Background:
    Homomorphic CNOT slide. Protocol and key update: the manuscript's
    Non-CCZ Evaluation section.
  }
\end{frame}
